\documentclass[10pt,a4paper]{amsart}
\usepackage[margin=19mm]{geometry}
\usepackage{amsmath,amssymb,mathtools}
\usepackage[scr=rsfs,cal=euler]{mathalfa}
\usepackage{xfrac}
\usepackage[unicode,hidelinks,bookmarks=false]{hyperref}
\newcommand{\ie}{i.e.\ }
\newcommand{\cf}{cf.\ }
\newcommand{\resp}{respectively\ }
\newcommand{\Subsection}{\subsection}
\providecommand{\nobreakdash}{\nobreak\hbox{-}}
\setlength{\emergencystretch}{2em}
\multlinegap0pt
\usepackage{bm}
\usepackage{bbm}
\usepackage{dsfont}
\usepackage{xspace}
\usepackage{cancel}
\usepackage{mathtools}
\usepackage{amsthm}
\makeatletter
\@ifundefined{theorem}{\newtheorem{theorem}{Theorem}}{}
\makeatother
\newcommand{\F}{\mathcal{F}}
\newcommand{\Ghat}{\widehat{G}}
\title[On continuous embeddings of quantum Sobolev spaces into Scha]{On continuous embeddings of quantum Sobolev spaces into Schatten classes $\mathfrak{H}\_\gamma^{s,p}(G,H) \hookrightarrow S\_p(H)$}
\author{Alexander Plakhotnikov}
\date{Source: arXiv:2509.24135v2 (CC BY 4.0)}
\begin{document}
\maketitle

\noindent\emph{Source and licence.} On continuous embeddings of quantum Sobolev spaces into Schatten classes $\mathfrak{H}_\gamma^{s,p}(G,H) \hookrightarrow S_p(H)$. Version arXiv:2509.24135v2 (CC BY 4.0), distributed under the Creative Commons Attribution 4.0 International licence, as recorded in the arXiv licence metadata for this exact version and shown on the abstract page https://arxiv.org/abs/2509.24135v2. Full rights notice: licenses/arxiv-2509.24135-CC-BY-4.0.txt.\par\smallskip

\noindent\textit{Editorial scope.} Counted unit: the Theorem whose source label is \texttt{thm:topological-iso} in the article (source0.tex lines 162--165, source label \texttt{thm:topological-iso}), delivered together with the article's own complete original proof of it (source0.tex lines 167--209, one \texttt{proof} environment of 43 lines). No mathematics is written, repaired or re-derived: statement and proof are the article's own text line for line. Required context retained uncounted: none. Every definition, display and label of those lines is retained unchanged. Omitted from this package and not redistributed: 1--161, 166--166, 210--557. Nothing inside a retained excerpt is omitted or altered: every retained body is delivered exactly as the article's lines are written, so no omission receipt of any kind accompanies this package. Citations: the retained excerpts carry no citation command at all, so no bibliography entry of the article is transcribed and no reference is redistributed. Reference adaptation: none was needed and none was made; every reference inside the delivered unit resolves inside it, so no unresolved reference or citation is produced. Printed slips kept literal: no printed word, symbol, hypothesis or number of the counted statement or of its proof was altered. Layout and class: the article's own class is \texttt{article} with packages including inputenc, fontenc, babel, amsmath, amsfonts, amssymb, amsthm, mathtools, geometry, tabularx, footmisc, xcolor, hyperref, comment; the wrapper is the lane's pinned \texttt{amsart} editorial wrapper at 10pt on A4, which restates the theorem-like environments the retained text needs and adds the article's own macro definitions that the retained text needs (\texttt{\textbackslash F}, \texttt{\textbackslash Ghat}), with extra wrapper packages (bm, bbm, dsfont, xspace, cancel, mathtools). The delivered numbering is the unit's own. Assets: no retained span contains a figure environment, a graphics-inclusion command, a PDF-page inclusion, a table or a TikZ picture; no image file is copied into this package, the package declares no asset dependency and no image file is redistributed with it. Licence: version arXiv:2509.24135v2 of this article is distributed under the Creative Commons Attribution 4.0 International licence, as recorded in the arXiv licence metadata for this exact version and shown on the abstract page https://arxiv.org/abs/2509.24135v2. A full rights notice is delivered at licenses/arxiv-2509.24135-CC-BY-4.0.txt. Characters: every retained excerpt is pure ASCII, so the delivered engine, XeLaTeX with its default Latin Modern text font, reports no missing character.
\begin{theorem}
\label{thm:topological-iso}
Let $1 < p < 2$. $\mathfrak{H}_\gamma^{s,p}(G, H)$ is a Banach space. The mapping \\$\Phi : T \mapsto (1+\gamma(\cdot)^2)^{s/2}\F_U(T)$ establishes a topological isomorphism between $\mathfrak{H}_\gamma^{s,p}(G, H)$ and its image in $L^q(\hat{G})$. The image of $\Phi$ is a closed subspace of $L^q(\hat{G})$, and $\Phi$ is a homeomorphism onto its image.
\end{theorem}

\begin{proof}

\medskip
1. Let $(T_n)$ be a Cauchy sequence in $\mathfrak{H}_\gamma^{s,p}(G,H)$. We have
\[
\|T_n - T_m\|_{\mathfrak{H}_\gamma^{s,p}}
  = \big\|(1+\gamma(\cdot)^2)^{s/2}\big(\F_U(T_n)-\F_U(T_m)\big)\big\|_{L^q(\widehat G)}
\]
Thus the condition that $(T_n)$ is Cauchy in $\mathfrak{H}_\gamma^{s,p}$ is 
\[
\forall\,\varepsilon>0\ \exists N\ \forall n,m\ge N:\ 
\big\|(1+\gamma(\cdot)^2)^{s/2}\big(\F_U(T_n)-\F_U(T_m)\big)\big\|_{L^q(\widehat G)}<\varepsilon
\]
In other words, the sequence
\[
\big((1+\gamma(\cdot)^2)^{s/2}\F_U(T_n)\big)_{n\in\mathbb N}
\]
is a Cauchy sequence in $L^q(\widehat G)$. Since $L^q(\widehat G)$ is complete, there exists
\[
g\in L^q(\widehat G)\quad\text{such that}\quad
(1+\gamma(\cdot)^2)^{s/2}\F_U(T_n)\to g \ \text{ in }L^q(\widehat G)
\]

\bigskip

\textsc{2: $\Phi$ is continuous.} This follows immediately from the definition: $\|\Phi(T)\|_{L^q} = \|T\|_{\mathfrak{H}_\gamma^{s,p}}$.

\bigskip

\textsc{$\Phi$ is injective.} If $\Phi(T) = 0$, then $(1 + \gamma(\cdot)^2)^{s/2} \F_U(T) = 0$ a.e. Since $\gamma(\xi) > 0$, we have $\F_U(T) = 0$ a.e. By injectivity of $\F_U$ on $S_p(H)$ (from Hausdorff-Young), we get $T = 0$.

\bigskip

\textsc{4.Image of $\Phi$ is closed in $L^q(\Ghat)$.}  
Let $(T_n)$ be a sequence in $\mathfrak{H}_\gamma^{s,p}(G,H)$ such that $\Phi(T_n)$ converges in $L^q(\hat{G})$ to some $g$. Since $(T_n)$ is bounded in the reflexive Banach space $S_p(H)$, by the Banach--Alaoglu theorem there exists a subsequence $(T_{n_k})$ weakly converging to some $T \in S_p(H)$. The continuity of $\Phi$ implies that $\Phi(T_{n_k})$ converges weakly to $\Phi(T)$ in $L^q(\hat{G})$. As the subsequence $\Phi(T_{n_k})$ also converges strongly to $g$, the weak limit must coincide with the strong limit, hence $\Phi(T) = g$. Therefore, $g$ lies in the image of $\Phi$, proving that the image is closed in $L^q(\hat{G})$.

Since $||T_n||_{\mathfrak{H}^{s,p}_\gamma} = ||f_n||_{L^q}$ is bounded, and $S_p(H)$ is reflexive  for $1 < p < \infty$, there exists a weakly convergent subsequence $(T_{n_k}) \rightarrow T \text{  in } S_p(H)$.


\bigskip

\textsc{5: $\Phi^{-1}$ is continuous on $\text{Im}(\Phi)$.} By the Open Mapping Theorem, since $\Phi: \mathfrak{H}_\gamma^{s,p}(G,H) \to \text{Im}(\Phi)$ is a continuous bijection between Banach spaces, it is a homeomorphism.
\end{proof}

\end{document}
